\documentclass[../main.tex]{subfiles} % 使用 subfiles 文档类，指定主文档
\graphicspath{{\subfix{../image/}}} % 指定图片目录，后续可以直接使用图片文件名。

% 例如:
% \begin{figure}[H]
% \centering
% \includegraphics[width=0.3\textwidth]{图.png}
% \caption{}
% \label{figure:图}
% \end{figure}
% 注意:上述\label{}一定要放在\caption{}之后,否则引用图片序号会只会显示??.

\begin{document}

\chapter{位势方程(Poisson 方程)}

\begin{definition}
设 $\Omega\subset \mathbb{R}^n$,我们称方程\eqref{eq:poisson}为\textbf{位势方程}或$\mathbf{Poisson}$\textbf{方程}.
\begin{align}
-\Delta u = f(x), \label{eq:poisson}
\end{align}
其中 $u=u(x)$，$x=(x_1,x_2,\cdots,x_n)\in\Omega\subset\mathbb{R}^n$，$f(x)$ 是一个已知函数。当非齐次项 $f\equiv 0$ 时，Poisson 方程 \eqref{eq:poisson} 简化为$\mathbf{Laplace}$\textbf{方程}
\begin{align}
\Delta u = 0. \label{eq:laplace}
\end{align}
为了求解 Poisson 方程 \eqref{eq:poisson}，通常还要提适当的边值条件。对于 Poisson 方程 \eqref{eq:poisson} 来说，这样的边值条件通常分为以下三类:
\begin{enumerate}[(1)]
\item \textbf{第一边值条件:} 已知函数 $u$ 在区域边界 $\partial\Omega$ 上的值，即
\begin{align*}
u(x) = g(x), \quad x\in\partial\Omega.
\end{align*}
这一类边值条件亦称为\textbf{ Dirichlet 边值条件}，相应的边值问题也称为\textbf{第一边值问题}或 \textbf{Dirichlet 边值问题}。

\item \textbf{第二边值条件:} 已知函数 $u$ 在区域边界 $\partial\Omega$ 上的法向导数，即
\begin{align*}
\frac{\partial u}{\partial n} = g(x), \quad x\in\partial\Omega,
\end{align*}
其中 $n$ 是 $\partial\Omega$ 的单位外法向量。这一类边值条件亦称为 \textbf{Neumann 边值条件}，相应的边值问题也称为\textbf{第二边值问题}或 \textbf{Neumann 边值问题}。

\item \textbf{第三边值条件:} 已知函数 $u$ 和其在区域边界 $\partial\Omega$ 上的法向导数的一个线性组合，即
\begin{align*}
\frac{\partial u}{\partial n} + \alpha(x)u(x) = g(x), \quad x\in\partial\Omega,
\end{align*}
其中 $n$ 是 $\partial\Omega$ 的单位外法向量，$\alpha(x)>0$。相应的边值问题也称为\textbf{第三边值问题}。
\end{enumerate}
\end{definition}

\subfile{Chapter-02/section-01.tex}

\subfile{Chapter-02/section-02.tex}

\subfile{Chapter-02/section-03.tex}

\subfile{Chapter-02/section-04.tex}

\subfile{Chapter-02/section-05.tex}

\subfile{Chapter-02/section-06.tex}

\subfile{Chapter-02/section-07.tex}

\subfile{Chapter-02/section-08.tex}

\subfile{Chapter-02/section-09.tex}

\subfile{Chapter-02/section-10.tex}

\end{document}